% !TeX root = ../navier-stokes-manuscript.tex
% ============================================================
% 1.2 Where the Problem Is Today
% ============================================================
\subsection{The Classical Estimates Before Closure}\label{sub:HistoricalEstimates}

The local theorem has launched a smooth fluid.  To keep it alive, we must know
which part of that smoothness survives for as long as the solution exists.
Kinetic energy survives globally, but at a derivative level too low to restart
the classical history.  The classical theory measures the exact distance
between those two facts.

\divider{Energy reaches the endpoint before the gradient does}

\begin{proposition}[Classical kinetic-energy identity]\label[proposition]{prop:ClassicalKineticEnergyIdentity}
Let \((u,p)\) be a sufficiently decaying smooth solution of
\eqref{eq:NavierStokesSystem}.  At every time in its interval of existence,
\begin{equation}\label{eq:AppendixKineticEnergyIdentity}
  \lVert u(t)\rVert_{L^2}^2
  +
  2\nu\int_0^t\lVert\nabla u(s)\rVert_{L^2}^2\,ds
  =
  \lVert u_0\rVert_{L^2}^2.
\end{equation}
\end{proposition}
\begin{proof}
Pair the momentum equation with \(u\) in space.  The transport term carries
energy through the fluid but cannot change its total amount:
\[
  \int_{\mathbb R^3}(u\cdot\nabla)u\cdot u\,dx
  =
  \frac12\int_{\mathbb R^3}u\cdot\nabla|u|^2\,dx
  =0.
\]
Pressure also does no total work, since
\[
  \int_{\mathbb R^3}\nabla p\cdot u\,dx
  =
  -\int_{\mathbb R^3}p\,\nabla\cdot u\,dx
  =0.
\]
The viscous term alone has a sign, and integration by parts gives
\[
  \nu\int_{\mathbb R^3}\Delta u\cdot u\,dx
  =
  -\nu\lVert\nabla u\rVert_{L^2}^2.
\]
Time integration is exactly \eqref{eq:AppendixKineticEnergyIdentity}.\qedhere
\end{proof}
Thus the velocity has a pointwise-in-time \(L^2\) ceiling, while the gradient is
paid only after its square has been accumulated over time.

Leray’s construction preserves precisely this information without assuming the
solution remains smooth.  Regularization and compactness produce a global weak
solution with \parencite{Leray1934}
\[
  u\in
  L^\infty(0,T;L^2(\mathbb R^3))
  \cap
  L^2(0,T;H^1(\mathbb R^3))
\]
on each finite interval.  The weak history is global, but this estimate by
itself neither makes arbitrary three-dimensional Leray solutions unique nor
returns the missing smoothness.

The spatial half of the energy class improves through
\[
  H^1(\mathbb R^3)\hookrightarrow L^6(\mathbb R^3)
\]
and therefore
\[
  u\in L^2(0,T;L^6(\mathbb R^3)).
\]
The nonlinear term consumes that gain immediately:
\[
  \lVert(u\cdot\nabla)u\rVert_{L^{3/2}}
  \leq
  \lVert u\rVert_{L^6}
  \lVert\nabla u\rVert_{L^2},
\]
so transport lies in \(L^1_tL^{3/2}_x\), enough for the weak equation to make
sense.  It does not bound \(\lVert\nabla u(t)\rVert_2\) at each time.  Taller
peaks can occupy shorter intervals and still leave the squared time integral
finite.

In three dimensions, that missing instantaneous gradient is exactly where a
vortex can amplify itself.  Put
\[
  \omega:=\nabla\times u.
\]
For the divergence-free decaying fields at hand,
\[
  \lVert\nabla u\rVert_{L^2}^2
  =
  \lVert\omega\rVert_{L^2}^2,
\]
so the time-integrated gradient payment is the time-integrated enstrophy.
Taking curl reveals its evolution:
\[
  \partial_t\omega+(u\cdot\nabla)\omega
  =
  (\omega\cdot\nabla)u+\nu\Delta\omega.
\]
The same velocity gradient that the energy estimate controls only in time now
stretches and tilts \(\omega\).  Pairing this equation with \(\omega\) gives
\[
  \frac12\frac{d}{dt}\lVert\omega\rVert_{L^2}^2
  +
  \nu\lVert\nabla\omega\rVert_{L^2}^2
  =
  \int_{\mathbb R^3}(\omega\cdot\nabla)u\cdot\omega\,dx.
\]
Nothing fixes the sign of the stretching integral, and kinetic energy does not
force it below dissipation at each instant.  The global estimate has therefore
reached the whole interval without reaching the derivative height required for
a classical restart.

\divider{The derivative height that buys another interval}

For an integer \(m\geq0\), collect the first \(m\) spatial derivative levels in
\[
  \lVert u\rVert_{H^m}^2
  :=
  \sum_{|\alpha|\leq m}
  \lVert\partial^\alpha u\rVert_{L^2}^2,
\]
and use the Fourier definition at fractional order.  In three dimensions,
Sobolev embedding gives \parencite[Theorem~7.28]{Cerda2010}
\[
  H^s(\mathbb R^3)\hookrightarrow C^k(\mathbb R^3)
  \qquad
  \text{when}
  \qquad
  s>k+\frac32.
\]
In particular,
\[
  s>\frac52
  \qquad\Longrightarrow\qquad
  H^s(\mathbb R^3)\hookrightarrow W^{1,\infty}(\mathbb R^3).
\]
Above \(5/2\), one Sobolev norm controls the pointwise gradient that multiplies
every differentiated energy row.

Suppose, then, that on \([0,T)\)
\[
  \sup_{0\leq t<T}\lVert u(t)\rVert_{H^s}
  \leq M
  \qquad
  \text{for some }s>\frac52.
\]
At any fixed finite height \(m\geq s\), differentiate the equation, pair each
row with the matching derivative of \(u\), and sum.  Incompressibility cancels
the top transport term, leaving
\[
  \frac12\frac{d}{dt}\lVert u\rVert_{H^m}^2
  +
  \nu\lVert\nabla u\rVert_{H^m}^2
  \leq
  C_m\lVert\nabla u\rVert_{L^\infty}
  \lVert u\rVert_{H^m}^2.
\]
The assumed \(H^s\) roof bounds the coefficient, so Gronwall carries this
chosen \(H^m\) norm to time \(T\).  Its constant may depend on \(m\): smoothness
requires each finite height, not one uniform constant for infinitely many
heights.

Once the velocity derivatives are present, pressure is forced by them.  Taking
divergence gives
\[
  -\Delta p
  =
  \sum_{i,j=1}^3\partial_i\partial_j(u_i u_j),
\]
with spatial decay fixing its additive constant.  The elliptic relation returns
the matching spatial derivatives of \(p\); the momentum equation then returns
\[
  \partial_tu
  =
  \nu\Delta u-(u\cdot\nabla)u-\nabla p
\]
Repeating pressure recovery and the momentum equation produces every fixed
mixed derivative.  One time-uniform Sobolev roof above \(5/2\) therefore
recovers the full smooth velocity--pressure pair.

That roof also provides time.  Local well-posedness in \(H^s\) has a lifespan
controlled by \(s\), \(\nu\), and the initial \(H^s\) size
\parencite{FujitaKato1964,Kato1984}.  Restart at a late slice
\(u(\cdot,t_0)\).  If all such slices remain below \(M\), their lifespans share
a positive lower bound; choosing \(t_0\) close enough to \(T\) crosses it.
Uniqueness joins the restart to the original history.  Consequently, a finite
maximal time can occur only if
\[
  \limsup_{t\uparrow T_*}
  \lVert u(t)\rVert_{H^s}
  =\infty
  \qquad
  \text{for every }s>\frac52.
\]
The unpaid distance is now explicit: the energy class gives time-integrated
\(H^1\), while continuation needs time-uniform control above \(H^{5/2}\).

\divider{The scale on which concentration can hide}

The Navier--Stokes scaling explains why energy cannot simply be upgraded to the
missing roof.  From any solution \((u,p)\), form
\[
  u_\lambda(x,t)
  :=
  \lambda u(\lambda x,\lambda^2t),
  \qquad
  p_\lambda(x,t)
  :=
  \lambda^2p(\lambda x,\lambda^2t)
\]
which solves the same equation after space and time are contracted together.
Any invariant norm sees the rescaled structure at the same size even while its
physical width tends to zero.

The Ladyzhenskaya--Prodi--Serrin line balances the strength of that spatial
concentration against the time for which it persists
\parencite{Prodi1959,Serrin1962,Serrin1963,Ladyzhenskaya1967}.  A solution
continues through \(T\) if
\[
  u\in L^q(0,T;L^p(\mathbb R^3)),
  \qquad
  \frac{2}{q}+\frac{3}{p}\leq1,
  \qquad
  p>3.
\]
At equality, spatial intensity and temporal duration change by inverse factors
under the Navier--Stokes scaling; the norm is critical.

The limiting spatial member is \(L^\infty_tL^3_x\), since
\[
  \lVert u_\lambda(t)\rVert_{L^3}
  =
  \lVert u(\lambda^2t)\rVert_{L^3}.
\]
Escauriaza, Seregin, and \v{S}ver{\'a}k showed that even this endpoint condition
\[
  u\in L^\infty(0,T;L^3(\mathbb R^3))
\]
continues the solution through \(T\)
\parencite{EscauriazaSereginSverak2003}.  Their argument rescales a first
singularity into an ancient limit and eliminates it by backward uniqueness.
Hence a finite last time must make the critical \(L^3\) norm escape.

The same threat has a local description.  Around
\(z_0=(x_0,t_0)\), use the parabolic cylinder
\[
  Q_r(z_0)
  :=
  B_r(x_0)\times(t_0-r^2,t_0).
\]
Caffarelli, Kohn, and Nirenberg proved that small scale-invariant velocity and
pressure content on \(Q_r(z_0)\) makes a smaller cylinder regular
\parencite{CaffarelliKohnNirenberg1982}.  A singular point must therefore retain
a definite amount of that content at every sufficiently small scale.  Their
covering argument makes the possible singular set small in parabolic Hausdorff
measure, but one surviving point is enough to destroy global smoothness.

For smooth compactly supported whole-space data, Seregin sharpens the endpoint
statement to actual divergence along the terminal approach
\parencite{Seregin2012}:
\[
  \lim_{t\uparrow T_*}
  \lVert u(t)\rVert_{L^3(\mathbb R^3)}
  =\infty.
\]
Tao’s quantitative refinement forces a positive power of a triple logarithm of
the inverse time remaining along a sequence of terminal times
\parencite{Tao2021Quantitative}.  This lower rate is far slower than any power
law; the classical results still permit much faster concentration.

\divider{The one estimate still absent}

If the maximal time were finite, kinetic energy would remain bounded while the
critical \(L^3\) norm and every continuation-grade Sobolev norm escaped.
Locally, some parabolic cylinders would retain nonzero scale-invariant content
at arbitrarily small radii.  These results describe what a singularity must do;
they do not give the terminal-uniform estimate that prevents it.

The mismatch with energy is visible in the two scaling laws
\[
  \lVert u_\lambda(t)\rVert_{L^3}
  =
  \lVert u(\lambda^2t)\rVert_{L^3},
  \qquad
  \lVert u_\lambda(t)\rVert_{L^2}^2
  =
  \lambda^{-1}
  \lVert u(\lambda^2t)\rVert_{L^2}^2.
\]
A shorter structure retains its critical \(L^3\) size while carrying less
\(L^2\) energy.  A global kinetic-energy ceiling therefore cannot become a
scale-independent ceiling on the critical concentration.

The same scaling shortens the clock.  Repeatedly halving the active width gives
parabolic durations proportional to
\[
  1,\ \frac14,\ \frac1{16},\ \ldots,
\]
whose sum is finite.  Each stage can retain critical size while occupying less
time, and none of the estimates above forbids that arithmetic.  This is all the
Zeno schedule asserts: it exposes an available timing, not a constructed
singular Navier--Stokes history.

Closure therefore asks for one terminal-uniform continuation-grade norm on the
actual maximal history.
